\section{Unit template: dua-template}
\label{sec:dua_template}
The previous section showed how the consistent environment provided by DUA is designed and built to support an already wide variety of hardware platforms and host system configurations. Recalling that the intended user of this framework is a researcher or a developer tasked with the design of a prototype for a distributed autonomous system, what is left to do is to provide them with a development environment based on the base units discussed before. In other words, since frameworks inherently define a standard way of developing and organizing software as shown by the other solutions presented in \Cref{sec:prb_related_work}, DUA has to provide a \emph{template} for the development of a unit. This is the purpose of the next core DUA module: \texttt{dua-template} \cite{masocco_dua-template_nodate}.

The \texttt{dua-template} is implemented as a GitHub template repository, \emph{i.e.} a special type of Git code repository hosted on GitHub \cite{github_inc_github_nodate} which exists only to be \emph{forked}, \emph{i.e.}, copied to generate new projects which start with its contents. It contains a set of files that make up a structural template for the development of software modules for distributed systems as DUA units. Once appropriately forked on GitHub, it can be configured as the starting point for a new project as hereby described. The main goal of this project is to offer a common, shared work environment that minimizes the time and effort required to set up prototyping and experimenting of new solutions, as well as to provide a common framework for teams that can be versatile enough to support any development workflow. It aims at achieving so by providing three main features: workspace organization for both development and deployment, hardware and operating system abstraction via the base units and their contents, modular structure for easy integration of multiple components into a single macro-project or distributed architecture.

A key component of a software development process is today a good Integrated Development Environment (IDE). During the development of this work, from its embryonic stages to its current version, the IDE of choice has been Visual Studio Code (VS Code) by Microsoft \cite{microsoft_visual_nodate}. Visual Studio Code is a free, open-source, cross-platform IDE that can be used to develop software in many different languages. It is particularly useful for this use case, where it can be used to develop code in C/C++, Python, and other languages, and it can also be used to manage Git repositories and Docker containers. It is also very customizable and can be extended with plugins that add support for many different languages and tools. As the following sections show, this template is configured to work with VS Code out of the box, but it is just a suggestion: it is possible to use other IDEs, or none at all, by simply managing differently the files and directories that are included in the template according to the specific needs of the project at hand.

The description of the functionalities offered by this template starts from the explanation of the project layout that it provides, hereby informally referred to as the \emph{filesystem layout} of a unit. Then, it can be explained how to set up a DUA unit on top of one or more base units, and how to eventually set up a bigger project by integrating multiple units. The last part of this section is dedicated to the discussion of a GitHub workflow that is used to keep units and projects hosted on GitHub up to date with respect to the main \texttt{dua-template} version.

\subsection{Unit filesystem layout}
\label{sec:dua_template_layout}
The filesystem layout of a DUA unit is designed to be as simple and intuitive as possible, to allow the user to quickly understand the structure of the project and to easily find the files that they are looking for. The layout is also designed to be consistent across all units, to allow the user to quickly understand the structure of a new unit even if they have never seen it before. The layout is also designed to be flexible, to allow the user to easily add new files and directories to the project as needed. However, a standard set of directories is provided to help the user organize their project in a way that is consistent with the DUA framework. The layout is also designed to be compatible with the structure of the base units, to allow the user to easily integrate their software modules with the base units and to quickly deploy them on the target platform. A visual inspection of the layout of \cite{masocco_dua-template_nodate}, which is also depicted in \Cref{fig:duatemplatelayout}, shows that it resembles the organization of a colcon workspace discussed in \Cref{sec:ros2_buildsys}: this is due to the fact that the template is designed to be used also to develop software based on \ROS, and the colcon workspace is the standard way to organize \ROS projects. Nonetheless, a set of countermeasures are put in place to avoid adding build artifacts and similar contents to version control: almost every directory among those listed below contains an appropriate \texttt{.gitignore} file that prevents known file types to be added to Git, and the user is free to modify them as they see fit to customize this behavior. The best way to understand the \texttt{dua-template} filesystem layout is to examine the purpose and contents of each directory it contains, keeping in mind that the actual usage of some files is clarified later on.
\begin{figure}[t!]
  \centering
  \begin{adjustbox}{max width=\linewidth}\begin{tikzpicture}[
    % Global settings
    node distance=.5cm,
    box/.style={
        draw,
        rectangle,
        minimum width=2.2cm,
        minimum height=0.8cm,
        align=center
    },
    >={Stealth[scale=1.0]},
    shorten >=1pt, shorten <=1pt
  ]
    % Root node
    \node[box] (root) {\texttt{workspace}};

    % Second level (horizontal layout)
    % Position each child relative to the previous or to root
    \node[box] (github) [below=1.0cm of root, xshift=-9.5cm]  {\texttt{.github}};
    \node[box] (vscode) [right=0.5cm of github]               {\texttt{.vscode}};
    \node[box] (bin)    [right=0.5cm of vscode]               {\texttt{bin}};
    \node[box] (config) [right=0.5cm of bin]                  {\texttt{config}};
    \node[box] (docker) [right=0.5cm of config]               {\texttt{docker}};
    \node[box] (logs)   [right=0.5cm of docker]               {\texttt{logs}};
    \node[box] (src)    [right=0.5cm of logs]                 {\texttt{src}};
    \node[box] (tools)  [right=0.5cm of src]                  {\texttt{tools}};

    % Third level: directories under "docker"
    \node[box] (dockerchild1) [below left=1.0cm and 0.5cm of docker]  {\texttt{container-x86-base}};
    \node[box] (dockerchild2) [below right=1.0cm and 0.5cm of docker] {\texttt{...}};

    % Connections from root to second level
    \draw[->] (root.south) -- (github.north);
    \draw[->] (root.south) -- (vscode.north);
    \draw[->] (root.south) -- (bin.north);
    \draw[->] (root.south) -- (config.north);
    \draw[->] (root.south) -- (docker.north);
    \draw[->] (root.south) -- (logs.north);
    \draw[->] (root.south) -- (src.north);
    \draw[->] (root.south) -- (tools.north);

    % Connections from "docker" to third-level children
    \draw[->] (docker.south) -- (dockerchild1.north);
    \draw[->] (docker.south) -- (dockerchild2.north);
  \end{tikzpicture}\end{adjustbox}
  \caption{Directory tree of a DUA unit based on \texttt{dua-template} \cite{masocco_dua-template_nodate}.}
  \label{fig:duatemplatelayout}
\end{figure}

\subsubsection{\texttt{.github}}
\label{sec:dua_template_layout_github}
This directory contains the GitHub workflow files that are used to keep the unit up to date with respect to the main \texttt{dua-template} repository. The workflow is based on GitHub Actions, a feature of GitHub that allows to run automated tasks when certain events occur in a repository. It is contained in the file \texttt{sync-dua-template.yaml} in the \texttt{workflows} subdirectory, and it is triggered when a new commit is pushed to the repository or at a specific time of the day. It is presented in \Cref{sec:dua_template_update}.

\subsubsection{\texttt{.vscode}}
\label{sec:dua_template_layout_vscode}
This directory contains configuration files for Visual Studio Code. They are present to provide each user with a good starting point for their development environment, and they can be customized as needed. Keep in mind that they are applicable since, as discussed later on, Docker containers based on DUA base units for development are also meant to be opened and managed with VS Code, and when this is done the settings for the IDE are automatically loaded from this directory. The files are:
\begin{itemize}
  \item \texttt{c\_cpp\_properties.json}: this file contains the configuration for the C/C++ extension for VS Code, which is used to develop C/C++ code.
  \item \texttt{launch.json}: this file contains configurations to run software modules from a VS Code instance.
  \item \texttt{settings.json}: this file contains the general settings for VS Code, such as the theme and the font size, as well as system paths to be used by the IDE.
  \item \texttt{tasks.json}: this file contains the tasks that can be run from VS Code, such as building the software modules or running tests.
\end{itemize}

\subsubsection{\texttt{bin}}
\label{sec:dua_template_layout_bin}
This directory is intended to contain any sort of standalone executables and scripts concerning a project. As for any other directory that has an intended purpose, the user is encouraged to put such files here, but as many other directories in the template, this is just a suggestion and the user is free to organize their project as they see fit. In the template, only the \texttt{dua\_setup.sh} script is present, which is used to set up a DUA-based project as explained in the following sections.

\subsubsection{\texttt{config}}
\label{sec:dua_template_layout_config}
This directory is intended to contain configuration files that concern and entire unit or project. Again, this is just a suggestion, and for the moment the template does not require any file to be present from the start in this directory.

\subsubsection{\texttt{docker}}
\label{sec:dua_template_layout_docker}
This directory starts empty, but will contain the configuration files for the Docker containers that are used to develop and run the software modules. As explained previously, there may be a Docker container for each base unit that a project is intended to run on, so this directory may have multiple subdirectories, one for each base unit. This aspect is also explained later on, since the contents of this directory are also set up by the \texttt{dua\_setup.sh} script.

\subsubsection{\texttt{logs}}
\label{sec:dua_template_layout_logs}
This directory is also empty at the beginning, because it is intended to contain log files generated by the software modules when they are run, including sampled data of experimental runs. To ensure that such large files are not added to the version control system by mistake, a \texttt{.gitignore} file is present in the directory that excludes all files in it.

\subsubsection{\texttt{src}}
\label{sec:dua_template_layout_src}
This is probably the most important and most used directory of the entire template. It is intended to contain the source code of the software modules that are developed for the unit, as well as subdirectories containing units that are added to a larger project. This directory is also meant to be found and managed by colcon, in case it is used as a build system for the project.

\subsubsection{\texttt{tools}}
\label{sec:dua_template_layout_tools}
This last directory also starts as an empty directory, and is also meant to contain source code and software modules which are under active development withing the context of the unit or project at hand. The difference with respect to the \texttt{src} directory is that the files in this directory are not meant to be built using colcon, which in fact ignores this directory thanks to the \texttt{COLCON\_IGNORE} file that is present in it.

\subsection{Configuring a new unit}
\label{sec:dua_template_configuration}
The setup of a new unit or project is a process that is meant to be as simple as possible, to allow the user to quickly start developing software modules for a distributed system. The setup process is performed by running the \texttt{dua\_setup.sh} script that is present in the \texttt{bin} directory of the template. The script is designed to be run from the root directory of the unit, in one of the ways shown in \Cref{lst:duasetup}.
\begin{center}
\begin{minipage}[t]{.9\textwidth}
\begin{lstlisting}[caption=Possible invocations of the \texttt{dua\_setup.sh} script., label={lst:duasetup}]
dua_setup.sh create [-a UNIT1,UNIT2,...] NAME BASE_UNIT
dua_setup.sh modify [-a UNIT1,UNIT2,...] [-r UNIT1,UNIT2,...] BASE_UNIT
dua_setup.sh delete BASE_UNIT
dua_setup.sh clear BASE_UNIT
\end{lstlisting}
\end{minipage}
\end{center}
An explanation of all the possible invocations of the script is given in the following sections. Before going further, another definition is in order here: with the term \emph{target} we identify a single Docker image that is used in a project, and that is built from a single Dockerfile pair plus many other configuration files discussed later on. A target can be used to run a single unit, or multiple ones, or even a whole control architecture; the name resembles the meaning of the term target in the context of a Makefile, and is used to refer to the same concept in this context: each unit or architecture can have multiple targets, \emph{i.e.}, Docker images, each one based on a different base unit. Thanks to base units, the software that makes up a unit or project is ensured to run appropriately and similarly across all supported targets.

\subsubsection{Creating a target}
\label{sec:dua_template_configuration_create}
The command at line 1 of \Cref{lst:duasetup} creates a new target, \emph{i.e.}, a Docker image for a unit or project, plus a set of configuration files. The \texttt{-a} option can be used to specify a list of units that will be integrated in the new target, a feature that is explained later on in \Cref{sec:dua_template_configuration_modify}, and the NAME argument is the name of the project. The \texttt{BASE\_UNIT} argument is the name of the base unit to use. A password is required for all projects since the container's internal user runs with elevated privileges and has full access to the host's network stack and hardware devices. It will be asked for during the creation of the target, and will be stored in hashed form in the target's Dockerfile.

\texttt{create} will create a new \texttt{container-BASE\_UNIT} directory inside \texttt{docker/}, and will copy and configure all the template files stored in \texttt{bin/dua-templates}. These include \texttt{Dockerfile} and \texttt{docker-compose.yaml} files, as well as many shell configuration scripts and other configuration files that are used to build and run the image. Once copied, they can be modified to fully support the specific features of the unit or project at hand. In particular:
\begin{itemize}
  \item \texttt{aliases.sh} is meant to contain shell aliases.
  \item \texttt{commands.sh} is meant to contain shell functions that execute commands that involve the software package that make up the unit or project at hand.
  \item \texttt{ros2.sh} contains configuration options and functions for the \ROS installation that is provided in every DUA image.
  \item Other files configure the shells and other basic programs like text editors.
\end{itemize}
Once a target has been created, it can be managed using Docker or VS Code extensions interchangeably. Containers are meant to be run in detached mode, \emph{i.e.}, with the \texttt{-d} option, and can be accessed using the \texttt{docker-compose exec} command as shown in \Cref{lst:duazsh}.
\begin{center}
\begin{minipage}[t]{.9\textwidth}
\begin{lstlisting}[caption=Docker Compose command to open an instance of the Zsh shell inside a DUA container., label={lst:duazsh}]
docker-compose exec NAME-BASE_UNIT zsh
\end{lstlisting}
\end{minipage}
\end{center}
This is due to how DUA containers are brought up, which is written at the end of the most important file among those generated by the setup script in this invocation: the \texttt{docker-compose.yaml} file. This file is meant to be used with Docker Compose (see \Cref{sec:docker}), either directly or through VS Code, to start the container specifying the many configuration options for the Docker Engine that would instead result in an unreasonably long command line for \texttt{docker run}. Each base unit requires specific configurations, which is why the setup script copies and configures the right files for the right base unit from the template ones stored in \texttt{bin/dua-templates}. Keep in mind that the contents of this file can be modified accordingly once the target configuration has been generated. The settings specified in the \texttt{docker-compose.yaml} file can be summarized as follows:
\begin{itemize}
  \item Build context and other useful options for the container build process.
  \item Environment variables that allow terminals and graphical applications to work properly.
  \item Host networking, shared memory and IPC namespace, to provide to applications running in the container full communication between the container and the host, as well as with remote hosts.
  \item Access to GPUs, if present.
  \item Access to all hardware devices with full privileges and \texttt{/dev} and \texttt{/sys} mounts.
  \item Access to the unit or project files from the host filesystem: the root directory of the unit is mounted in the container at \texttt{$\sim$/workspace}, and all the configuration files for the container respect this convention.
\end{itemize}
These settings, apart from those exclusively dedicated to the image build, make up the core of the integration of Docker in DUA as a tool to build and deploy replicable system environments: applications running inside a container configured in such a way have access to the host system in pretty much the same way they would if they were running directly on it, with direct access to any hardware device and driver ranging from the network stack to communication interfaces and even GPUs, if they are present. The advantage is that the entire installation and dependency configuration is performed inside a Docker image, avoiding the risk of tainting the host configuration with the risk of compromising it, leading to a reinstall that would be a timely process especially on SBCs. Furthermore, a Docker image can be simply packed and downloaded to any number of hosts once it is finalized, making the replication of a system environment a matter of seconds, compared to traditional solutions for, \emph{e.g.}, SBCs, that would involve flashing the board with a new OS image and then configuring it from scratch, manually or using scripts. Moreover, a Docker image constitutes a snapshot of the system environment at a given time, including the version of the software packages and libraries that are installed in it, which can be useful to reproduce experiments or to ensure that a system is always running the same software version, guaranteeing stability and reproducibility. The only overhead is due to the size of the images, as in \Cref{tab:dua_sizes}, but considering that once the containers are started there is no more overhead, due to the fact that the unit files are mounted from the host filesystem, the advantages of using Docker in this context are clear. Nonetheless, these options deactivate by default many security features that Docker enforces, to achieve full isolation of a process running inside a container with respect to the rest of the host system. Note that this work, as similar solutions that are starting to leverage Docker in the field of autonomous systems, some of which are discussed in \Cref{sec:prb_related_work}, does not intend to offer that kind of isolation for two main reasons. First, the end goal is that of easily replicating system configurations across multiple hosts, allowing the possibility of applications accessing the full set of resources offered by the host system: this implies that only a subset of the functionalities of Docker are necessary by the definition of the original problem. Moreover, the context of interest is that of a development environment for research and prototyping in which things have to be made to work and tested quickly first, and then eventually robustified and gradually isolated. This is why the default settings of the Docker containers are set to allow full access to the host system, but nonetheless they are fully exposed to the user of the framework, who can modify them as they see fit to achieve a configuration in which containers are properly isolated, being exposed only to a subset of the features and resources of the host system that are necessary for their applications to run. Some of these options have not been considered for simplicity, such as configuration of the \emph{secure computing mode} (\texttt{seccomp}) module to restrict access to specific system calls, but given the complete set of configuration files offered by the template, the user can easily add them to the container configuration.

Finally, an item that deserves a thorough explanation is the \texttt{command} item, which specifies the command that is run when the container is started. This is set by default to what is shown in \Cref{lst:duacmd}, which is essentially a Bash shell process that sleeps indefinitely, being woken up only by a \texttt{TERM} signal that triggers a graceful exit. This is done to keep the container running and to allow the user to access it at any time, and in any way they see fit, \emph{i.e.}, opening within it multiple shells of the flavor they prefer, or running graphical applications, or running applications, or accessing server applications running within it from remote hosts, or connecting to it using a VS Code instance, etc. Keep in mind that these containers are meant to be deployed on autonomous systems, so they should stay active with an arbitrary set of processes running within them. This is again not the originally intended use of Docker, \emph{i.e.}, that of running single microservices inside each container, but is instead a new way of using Docker to provide a complete, versatile, and consistent environment, which is currently being exploited also by other frameworks in different contexts.
\begin{center}
\begin{minipage}[t]{.9\textwidth}
\begin{lstlisting}[caption=Default command of a DUA container: an idle shell process., label={lst:duacmd}]
/bin/bash -c "trap 'exit 0' TERM; sleep infinity & wait"
\end{lstlisting}
\end{minipage}
\end{center}

\subsubsection{Modifying a target and integrating units}
\label{sec:dua_template_configuration_modify}
The remaining core functionalities of \texttt{dua-template} concern the configuration of a unit to run the software modules that it contains. In this context, the main aim of DUA is not only to provide a stable and reproducible way to do so, but also a mechanism to simply integrate a unit alongside multiple other ones in a bigger project, in a modular and simple fashion.

Since units are software modules in essence, the system configurations they require usually boil down to two categories: system settings, like environment variables and configuration files, and dependencies, such as packages to install. Recalling how base units are set up as explained in \Cref{sec:dua_foundation}, these steps should be somehow added to the image build process on top of what the base unit already provides. This is the purpose of the \texttt{Dockerfile} file, generated by the \texttt{dua\_setup.sh} script and placed in the directory of every configured target.

The interested reader can find templates for the \texttt{Dockerfile} in \cite{masocco_dua-template_nodate}. As for the previous discussion, this section aims to guide the inspection of such templates by detailing their contents. The operations performed by the \texttt{Dockerfile} are, in order:
\begin{enumerate}
  \item Start from the target's base unit image.
  \item Perform unit-specific operations.
  \item Configure the internal user of the container.
  \item Create mount points for configuration files and the workspace.
  \item Configure the shell.
\end{enumerate}
While most steps are standard and should not be altered in general, although the user is given complete control over how the framework works, the second is the one where manual intervention is encouraged. This is because, initially, that step is nonexistent, as illustrated by \Cref{lst:imgsetup}.
\begin{center}
\begin{minipage}[t]{.9\textwidth}
\begin{lstlisting}[caption=Unit configuration step in a Dockerfile as generated by the template., label={lst:imgsetup}]
# NAME SETUP START #
# NAME SETUP END #
\end{lstlisting}
\end{minipage}
\end{center}
\texttt{NAME} is, of course, replaced by the name given to the unit by the user as in line 1 of \Cref{lst:duasetup}. Between those two lines, the user can add any Dockerfile directive they require to configure the system environment of the container to support the software that will make up the unit. Once built, the container will have all the features and functionalities of its base unit, plus what the user added to support their own software module.

This section of the \texttt{Dockerfile}, and the way it is delimited, are instrumental to achieve the desired level of modularity, making the integration of multiple units in a single project again a matter of a few commands, listed in the remaining lines of \Cref{lst:duasetup}.

Suppose that a unit must be integrated in a project. The first step consists in adding the unit codebase to the one of the project, \emph{i.e.}, copying it in the \texttt{src/} directory. It is assumed that all codebases are managed by the Git version control system, as stated in \Cref{sec:dua_assumptions}, which offers two functionalities meant to this purpose: submodules and subtrees, illustrated in \Cref{sec:git}. They are both capable of linking the unit's repository to the project one, but they are meant for different use cases.

Submodules are the preferred way to include units in a DUA project when the user plans to actively work on the units themselves, while working on the project. A submodule is a Git repository that is included in another Git repository as a subdirectory. This means that the submodule is a separate repository, with its own history, branches, and tags, that is included in the parent repository as a subdirectory. This is particularly useful when the submodule is actively developed, and when changes to it must be pushed back to the original repository. The main drawback of submodules is that they require additional commands to be managed, and that special attention must be paid to the point of the Git history in which the submodule is found before making any changes inside it, or worse, committing them. Moreover, the reference commit of the submodule in the parent repository must always be updated manually. Nonetheless, submodules are a very powerful tool that, when used correctly, can make the simultaneous development of multiple projects much easier. For this reason, given the degree of modularity that DUA aims to achieve, they are supported by \texttt{dua-template} and their usage is encouraged as the preferred way to include units in a DUA project in any case.

Subtrees, on the other hand, are the preferred way to include units in a DUA project when the user does not plan to actively work on the units themselves. Essentially, integrating another Git repository as a subtree still ensures the possibility of specifying the commit, branch, or tag of the repository to be integrated, but requiring a simpler management with respect to submodules, since files and Git histories will be merged in those of the parent repository, with little to no additional commands required to keep them in sync, resulting in a single repository that contains all the files and histories of its subtrees. However, they have one major drawback: they may make Git history management more computationally difficult, since the history of the subtree is merged in the one of the parent repository. This often leads to Git having to go through the whole history of the subtree to find the changes that are relevant to the parent repository, to ensure that a coherent history is maintained; this could easily become a slow process, which is why the use of subtrees is discouraged when the subtree is actively developed.

Then, as for every DUA codebase, a target for the project must be configured. What follows can be applied to any number of targets a project shall support. The \texttt{modify} command of the \texttt{dua\_setup.sh} script, at line 2 of \Cref{lst:duasetup}, is intended for this purpose. The \texttt{-a} option is used to integrate one or more units into the specified target, while the dual \texttt{-r} option is used to remove one or more units from the target. To ensure that all the system configurations that the supplied units require are performed at the same time within that target, what the script does in case of an addition is simply to look for the relevant units in the \texttt{src} directory, and then for their corresponding target directory; in case the target at hand is, \emph{e.g.}, \texttt{x86-base} and the unit is called \texttt{sensor-driver}, the script will look for the directory \texttt{src/sensor-driver/} \texttt{docker/container-x86-base}. It the search is successful, the script will proceed to copy the unit configuration section of the \texttt{Dockerfile}, delimited as in \Cref{lst:imgsetup}, within the target's \texttt{Dockerfile} and into its own unit configuration section. This way, the user can be sure that all the system configurations that are required to run the software module are performed at the same time, and in the same way, for all the units that are integrated in the target. The same operation is performed in case of a removal, but in reverse: the script will look for the unit configuration section in the target's \texttt{Dockerfile}, and will remove it if found. If some units need to be set up in a specific order, it is up to the user to specify them in the desired order in the \texttt{-a} option of the script. This description also clarifies the remaining reasons behind the filesystem layout of a unit, and the presence of the \texttt{Dockerfile} in the subdirectories of the \texttt{docker/} directory, as illustrated previously in \Cref{sec:dua_template_layout}.

The commands at line 3 and 4 of \Cref{lst:duasetup} are pretty straightforward and are provided for convenience. The \texttt{delete} command can be used to quickly and easily remove a target from a unit or project, effectively deleting its subdirectory from the \texttt{docker/} directory. The \texttt{clear} command can be used to remove all the lines between the delimiters of the unit configuration section of a target's \texttt{Dockerfile}, effectively resetting the target to its base unit configuration.

\subsection{Updating dua-template in units and projects}
\label{sec:dua_template_update}
As every software component, and because the DUA project is under active development, \texttt{dua-template} is subject to updates. Such updates can be of various types, ranging from bug fixes to new features, and they are meant to be applied to all units and projects that are based on the template. To ensure that all units and projects are up to date with respect to the main \texttt{dua-template} repository, a GitHub workflow is provided in the \texttt{.github} directory of the template, as explained in \Cref{sec:dua_template_layout_github}. Every GitHub repository forked from the main \texttt{dua-template} repository will have this workflow enabled by default, and it will be triggered when a new commit is pushed to the \texttt{master} branch of such repository or at a specific time of the day. As for any feature of DUA, the user can modify or disable this workflow as they see fit.

The interested reader can find the code of the workflow at either \cite{masocco_dua-template_nodate} or in \Cref{lst:duawf} in \Cref{app:workflow}. The workflow, meant to run on virtual machines called \emph{runners} provided by GitHub, clones both the main \texttt{dua-template} repository and the repository where it is running, and then it compares the two repositories to find out if the main repository has been updated, \emph{i.e.}, either modified or created or even deleted, with respect to the forked one. The files that the workflow compares are only those that are part of \texttt{dua-template}, \emph{i.e.}, it does not compare nor modify files belonging to targets that have already been created and configured. This behavior is intended by design, and due to two reasons: first, it is not possible to deterministically distinguish user-applied modifications to target files from those that are part of the template, and second, the user always has the last word on how their project is configured, having the ability to modify every template file as they see fit to ensure that their software works correctly. Thus, the user is the only one that can decide whether and how to apply the update, either manually reconfiguring or regenerating the targets as explained in \Cref{sec:dua_template_configuration} using the updated template files. To provide a convenient way of doing so, the workflow creates a pull request in the forked repository with the changes to the template portion of the codebase, and the user can then review it, update their unit files, and merge the pull request to apply the changes.
